\documentclass[11pt]{article}

% ---------------------------------------------------------------------------
% WORKSHOP / arXiv preprint draft -- SKELETON ONLY (prose to be authored).
% Swap \documentclass for the venue style file once a target is chosen.
% Story (LOCKED 2026-06-08):
%   Headline   : WE WIN -- the cheapest spectral-polar member is the fastest
%                LoRA optimizer.
%   Protagonist: diagonal Shampoo+polar.
%   Claim      : it matches the expensive members and beats AdamW 1.5-1.7x
%                across 14 (model,dataset,rank) settings -- the expensive
%                machinery is dead weight.
% ---------------------------------------------------------------------------

\usepackage{amssymb,amsmath,amsthm,amsbsy,amsfonts}
\usepackage{mathtools}
\usepackage{bm}
\usepackage{graphicx}
\usepackage{booktabs}
\usepackage{algorithm}
\usepackage{algpseudocode}
% Looser line spacing inside algorithm blocks for readability.
\makeatletter
\newcommand{\algspacing}{\setlength{\@tempdima}{0pt}\renewcommand{\baselinestretch}{1.25}\selectfont}
\makeatother
\algrenewcommand\algorithmicrequire{\textbf{Input:}}
% Phase header: blank line + vertical space above, bold label.
\algnewcommand\algorithmicphase{\vspace{5pt}\textbf}
\algnewcommand\Phase{\Statex\algorithmicphase}
\usepackage{enumitem}
\usepackage{xcolor}
\usepackage[colorlinks,citecolor=blue,urlcolor=blue]{hyperref}
\usepackage[nameinlink]{cleveref}
\usepackage[firstinits=true,backend=bibtex,style=numeric-comp,citestyle=numeric-comp,sorting=nyt,url=false,arxiv=false,isbn=false,doi=false]{biblatex}
\bibliography{references}

\graphicspath{{figs/}}

% --- macros -----------------------------------------------------------------
\renewcommand{\vec}[1]{\bm{#1}}
\newcommand{\mat}[1]{\bm{#1}}
\newcommand{\R}{\mathbb{R}}
\newcommand{\E}{\mathbb{E}}
\DeclareMathOperator{\diag}{diag}
\DeclareMathOperator*{\argmin}{arg\,min}
\DeclareMathOperator*{\argmax}{arg\,max}
\newcommand{\smax}{\sigma_{\max}}
\newcommand{\DW}{\Delta W}
\newcommand{\polar}{\varphi}                 % polar map  U\Sigma V^T -> U V^T
% Working method name -- rename in one place once finalized.
\newcommand{\methodname}{\textsc{Polar-LoRA}\xspace}
\usepackage{xspace}
% Visible marker for numbers pending the full-polar confirm/negate sweep.
\newcommand{\PLACEHOLDER}[1]{\textcolor{red}{[#1]}}
% theorem environments
\theoremstyle{plain}
\newtheorem{lemma}{Lemma}
\newtheorem{proposition}{Proposition}
\newtheorem{corollary}{Corollary}
\theoremstyle{definition}
\newtheorem{definition}{Definition}
\theoremstyle{remark}
\newtheorem{remark}{Remark}

\title{\methodname}  % working title -- revisit before submission
\author{}
\date{}

\begin{document}
\maketitle

% ===========================================================================
\begin{abstract}
% SKELETON: Farquhar 5-sentence formula -- one stub per sentence, author expands.
% Numbers are placeholders; RE-PULL from the loader before submission.
% (1) what:        introduce \methodname (Shampoo-style two-sided curvature whiten
%                  -> full polar -> operator-norm step radius).
% (2) why-hard:    spectral updates on the LoRA factors fail (factorization-dependent,
%                  factor norms diverge); known fixes appear in isolation, bundled
%                  into heavy methods -- unclear which ingredient carries the gain.
% (3) how:         keep the cheap essential parts -- diagonal two-sided whitener,
%                  full polar via PolarExpress, operator-norm step radius.
% (4) evidence:    14 (model,dataset,rank) settings, 4 base models, 4 corpora;
%                  reaches AdamW's loss in \PLACEHOLDER{1.2--1.7}$\times$ fewer steps;
%                  ablations isolate the polar cap; diagonal curvature suffices.
% (5) number:      near-free -- $\le 6\%$ wall per step at production batch.
\end{abstract}

% ===========================================================================
\section{Introduction}
\label{sec:intro}
% Hard limits: <= 1.5 pages; methods by page 2-3; 2-4 contribution bullets <=2 lines.
%
% Citation keys (verify entries in references.bib against arXiv before submission):
%   kang_specgf=2602.06385, imuon=2605.09238, spectron=2602.12429,
%   soap=2409.11321, adaprelora=2605.08734, mua=2602.06204,
%   polarexpress=2505.16932, muon=<verify>.

% SKELETON -- author writes prose. Beats to hit, in order:
% - Hook: LoRA is the default; the optimizer that trains the adapter is the
%   under-examined lever (practice defaults to AdamW).
% - Challenge: spectral/steepest-descent (Muon family) and Shampoo curvature both
%   reach target loss in fewer steps on full-weight training, but carrying them to
%   the LoRA factors is not straightforward -- per-factor polar is not invariant
%   to the dW = BA reparameterization; factor norms can diverge \cite{kang_specgf, imuon}.
% - Gap (remedies are partial, each bundled with expensive machinery):
%     * iMuon \cite{imuon}: reparam-invariance via partner-Gram-whitened polar,
%       but controls only the PRODUCT step -- per-factor step unbounded, no
%       lr transfer across rank.
%     * Spectron \cite{spectron}: per-factor operator-norm step radius, but for
%       native low-rank PRETRAINING, no rank-transfer claim.
%     * Shampoo/SOAP, AdaPreLoRA \cite{soap, adaprelora}: two-sided curvature, no
%       spectral cap; AdaPreLoRA cost scales with the dense weight.
%     * lr-scaling rules for LoRA \cite{mua}: rate-rank relationship, but for
%       sign/Adam-style updates, not spectral.
% - Approach: assemble the cheap union, ablate the expensive parts away.
%   \methodname = momentum -> two-sided DIAGONAL curvature whiten -> FULL polar
%   (spectral cap) -> operator-norm step radius rho = eta/(smax(A)+smax(B)).
%
% CONSTRAINT (author feedback): the lr-transfer claim is EMPIRICAL ONLY --
%   "our optimal learning rate transfers across rank; AdamW's does not."
%   Do NOT assert r^{-1/2}, do NOT say "theory predicts", do NOT claim what
%   \cite{mua} predicts for our method (mua may be cited as related work on
%   rank-aware lr, but state no prediction).

\paragraph{Contributions.}
\begin{itemize}[leftmargin=1.5em,itemsep=2pt,topsep=2pt]
  \item \textbf{Fast, cheap LoRA optimizer.} Reaches AdamW's eval loss in
        \PLACEHOLDER{1.2--1.7}$\times$ fewer steps across 14 settings; $\le 6\%$
        wall per step.
  \item \textbf{Operator-norm lr transfer (empirical).} Optimal rate stable across
        rank; AdamW's shrinks with rank; iMuon (no radius) does not transfer.
  \item \textbf{What carries the gain.} Ablation isolates the polar cap; cheap
        diagonal curvature suffices.
  \item \textbf{Benchmark.} A 14-setting LoRA-optimizer comparison under a
        speedup-to-AdamW metric over 4 base models and 4 corpora.
\end{itemize}

% Results-preview stub (author expands): speedup is task-dependent -- largest
% out-of-distribution and at higher rank
% (\PLACEHOLDER{$1.20\times$ in-distribution code vs.\ $1.57\times$ OOD, same model/rank});
% per-step overhead leaves the step-count speedup intact in wall-clock.

% ===========================================================================
\section{Setup and metric}
\label{sec:setup}
% SKELETON -- TODO author:
% - LoRA factor convention (A:(r,d_in), B:(d_out,r), dW=(alpha/r)BA).
% - Held-everything-fixed protocol (model, data+split, r/alpha, targets, seq len,
%   batch, dtype, constant LR no warmup).
% - Speedup-vs-AdamW metric (horizon / steps-to-match-AdamW-final-loss).
% - AdamW noise floor sigma as the unit for Delta claims.

% ===========================================================================
\section{Method: \methodname}
\label{sec:method}
% MAIN STORY (outlined): all three updates are the SAME object -- the linear-
% minimization oracle (LMO) of the gradient under a norm -- and differ ONLY in the
% norm / metric. 3.1 decoupled per-factor spectral norm; 3.2 coupled product/manifold
% metric; 3.3 learned two-sided curvature metric + explicit operator-norm budget.
% Prose is terse here by design (author expands); the equations are the content.
% Citation keys (verify in references.bib): kang_specgf, imuon, spectron,
% polarexpress, soap, mua, tilde_csd (compositional steepest descent -- no arXiv).

% --- setup + the steepest-descent (LMO) view (MAIN STORY) ---------------------
% All three updates are the same step (steepest descent under a norm) under
% different norms. Statements in body; proofs in \Cref{app:proofs}.
\paragraph{Setup.}
\begin{itemize}[leftmargin=1.5em,itemsep=1pt,topsep=2pt]
  \item LoRA factors $A\in\R^{r\times d_{\mathrm{in}}}$, $B\in\R^{d_{\mathrm{out}}\times r}$; adapter $\DW=(\alpha/r)\,BA$.
  \item $G\triangleq\partial L/\partial(BA)$: gradient of the loss w.r.t.\ the \emph{unscaled} product $BA$ (the $\alpha/r$ scale is absorbed into the learning rate $\eta>0$). Per-factor gradients $G_A=B^\top G$, $G_B=GA^\top$.
  \item $\dot A,\dot B$: the updates applied to the factors ($A\leftarrow A+\dot A$).
  \item $\polar(M)=UV^\top$ for the SVD $M=U\Sigma V^\top$ (sets each nonzero singular value to $1$); spectral norm $\|M\|_2=\smax(M)$; $\langle X,Y\rangle=\operatorname{tr}(X^\top Y)$; Gram matrices $S_A=AA^\top$, $S_B=B^\top B$.
\end{itemize}
\textbf{Steepest descent under a norm} (the LMO), for a factor gradient $M$ and a norm $\|\cdot\|$, is the update
\begin{equation}
  \eta\,\argmin_{\|U\|\le1}\langle M,U\rangle .
  \label{eq:lmo}
\end{equation}
The three variants specialize $\|\cdot\|$ to three different norms; \Cref{lem:lmo} is the case of the spectral norm $\|\cdot\|_2$.

\subsection{Naive Muon on the factors}
\label{sec:method-naive}
\paragraph{Norm.} The spectral norm $\|\cdot\|_2$, applied to each factor separately.
\begin{lemma}[Spectral-norm LMO]\label{lem:lmo}
$-\polar(M)\in\argmin_{\|U\|_2\le1}\langle M,U\rangle$ for any matrix $M$ (uniquely so when $M$ has full rank with distinct singular values). \emph{(\Cref{app:proofs}.)}
\end{lemma}
\noindent Applying \Cref{lem:lmo} to each factor gradient:
\begin{equation}
  \dot A=-\eta\,\polar(G_A),\qquad \dot B=-\eta\,\polar(G_B).
  \label{eq:naive}
\end{equation}
Two defects on a factored weight:
\begin{itemize}[leftmargin=1.5em,itemsep=1pt,topsep=2pt]
  \item \emph{Factorization-dependent.} The product $\DW=BA$ is invariant under $(A,B)\mapsto(NA,BN^{-1})$ for invertible $N$, but the factor gradient transforms as $G_A\mapsto N^{-\top}G_A$. Since $\polar$ is not equivariant, \eqref{eq:naive} depends on the factorization, not on $\DW$ alone \cite{imuon}.
  \item \emph{Unbounded merged step.} The merged step
        $\|B\dot A+\dot B A\|_2\le\eta\big(\|A\|_2+\|B\|_2\big)$
        grows with the factor norms, which can diverge under the corresponding flow \cite{kang_specgf}.
\end{itemize}

\subsection{iMuon / compositional Muon}
\label{sec:method-imuon}
\paragraph{LMO (half-split).} Budget each factor's contribution to the product step separately --- the tractable surrogate \cite{tilde_csd} for the harder joint constraint $\|\dot B A+B\dot A\|_2\le\eta$:
\begin{equation}
  \min_{\dot A}\ \langle G_A,\dot A\rangle \ \ \text{s.t.}\ \ \|B\dot A\|_2\le\eta,
  \qquad
  \min_{\dot B}\ \langle G_B,\dot B\rangle \ \ \text{s.t.}\ \ \|\dot B A\|_2\le\eta .
  \label{eq:imuon-lmo}
\end{equation}
\begin{proposition}[Partner-Gram step]\label{prop:imuon}
The solution of \eqref{eq:imuon-lmo} whitens each factor by its partner's Gram root:
\begin{equation*}
  \dot A=-\eta\,S_B^{-1/2}\,\polar\!\big(S_B^{-1/2}G_A\big),\qquad
  \dot B=-\eta\,\polar\!\big(G_B S_A^{-1/2}\big)\,S_A^{-1/2}.
\end{equation*}
\emph{(\Cref{app:proofs}; iMuon \cite{imuon}, and the half-split of compositional Muon \cite{tilde_csd}.)}
\end{proposition}
\begin{corollary}[Product-step invariance]\label{cor:imuon}
$\|B\dot A\|_2=\eta$, independent of the factor norms --- so no runtime rescale is needed. \emph{(\Cref{app:proofs}.)}
\end{corollary}
\begin{itemize}[leftmargin=1.5em,itemsep=1pt,topsep=2pt]
  \item \emph{Same operator, other compositions.} Compositional Muon applies it to the attention products $W_QW_K^\top$ and $W_OW_V$ \cite{tilde_csd}.
  \item \emph{The gap we close.} The raw factor step is uncontrolled --- $\|\dot A\|_2$ scales as $\eta/\sigma_{\min}(B)$ --- and there is no per-factor operator-norm radius, so the learning rate does not transfer across rank.
\end{itemize}

\subsection{Ours: \methodname}
\label{sec:method-ours}
\paragraph{Motivation.} Naive Muon (\Cref{lem:lmo}) and iMuon (\Cref{prop:imuon}) are the same LMO under the Euclidean inner product, differing only in what they cap. \methodname keeps the LMO and changes the \emph{inner product} to a curvature-weighted one.

\paragraph{Curvature inner product.} Equip weight space with
\begin{equation}
  \langle X,Y\rangle_{P,Q}=\operatorname{tr}(X^\top P\,Y\,Q),\qquad
  \|X\|_{P,Q}=\|P^{1/2}XQ^{1/2}\|_2,
  \label{eq:ours-metric}
\end{equation}
where $P,Q\succ0$ are two diagonal curvatures --- per-coordinate gradient energy on the output and input space (EMAs, decay $\beta_2$),
\begin{equation*}
  P_{ii}=\mathrm{EMA}\,\big\|(G_B)_{i,:}\big\|_2^2,\qquad
  Q_{jj}=\mathrm{EMA}\,\big\|(G_A)_{:,j}\big\|_2^2,
\end{equation*}
and $\|\cdot\|_{P,Q}$ is the operator norm this inner product induces.

\paragraph{Derivation.} Cap the product step in this inner product and maximize alignment with the gradient:
\begin{equation}
  \max_{\dot A}\ \langle B\dot A,\,-G\rangle
  \quad\text{s.t.}\quad
  \|B\dot A\|_{P,Q}\le\tau .
  \label{eq:ours-opt}
\end{equation}
This is the LMO of \Cref{lem:lmo} in the curvature inner product. Whitening both sides reduces it to the spectral case, giving a closed form.
\begin{proposition}[Curvature-metric LMO]\label{prop:ours}
For $B$ of full column rank, the solution of \eqref{eq:ours-opt} is
\begin{equation*}
  \dot A = -\tau\,C_A^{-1/2}\,\polar\!\big(C_A^{-1/2}\,g_A\,Q^{-1/2}\big)\,Q^{-1/2},
  \qquad C_A=B^\top P B .
\end{equation*}
\end{proposition}
\medskip
\noindent \methodname applies \Cref{prop:ours} to the bias-corrected momentum $\hat M_A$:
\begin{equation}
  W_A \;=\; C_A^{-1/2}\,\polar\!\big(C_A^{-1/2}\,\hat M_A\,Q^{-1/2}\big)\,Q^{-1/2},
  \label{eq:ours-lmo}
\end{equation}
with $W_B$ the mirror ($C_B=AQA^\top$). \Cref{alg:ours} sets the step magnitude by the operator-norm budget $\rho=\eta/(\smax(A)+\smax(B))$, which caps the merged step independent of rank:
\begin{equation}
  \|B\dot A+\dot B A\|_2\le\rho\big(\smax(A)+\smax(B)\big)=\eta ,
  \label{eq:ours-prodcap}
\end{equation}
so the learning rate transfers across rank.

\begin{remark}[Relation to AdaPreLoRA]
\Cref{prop:ours} caps $Y_A$ in operator norm. Capping it in Frobenius norm instead, $\|Y_A\|_F\le\tau$, replaces the polar with the linear solve $\dot A\propto C_A^{-1}g_A Q^{-1}$ --- this is AdaPreLoRA \cite{adaprelora}, the same projection without the spectral cap.
\end{remark}

\begin{algorithm}[t]
\caption{\methodname{} step for one LoRA pair $(A,B)$ at iteration $t$}
\label{alg:ours}
\begin{algorithmic}[1]
\Require gradients $G_A,G_B$; momenta $M_A,M_B$; diagonal curvatures $P,Q$; decays $\beta_1,\beta_2$; rate $\eta$

\Phase{Momentum} (Nesterov lookahead, bias-corrected)
\State $M_A\gets\beta_1 M_A+(1-\beta_1)\,G_A$, \qquad $\hat M_A\gets\big(\beta_1 M_A+(1-\beta_1)G_A\big)/(1-\beta_1^{\,t})$
\State $M_B\gets\beta_1 M_B+(1-\beta_1)\,G_B$, \qquad $\hat M_B\gets\big(\beta_1 M_B+(1-\beta_1)G_B\big)/(1-\beta_1^{\,t})$
\Statex \hspace{1.2em}\emph{(EMA analog of Muon's $\hat g+\beta\,\mathrm{buf}$; plain $\hat M_A\!=\!M_A/(1-\beta_1^t)$ is the $-$Nesterov ablation.)}

\Phase{Diagonal curvature} (per-coordinate gradient energy)
\State $Q_{jj}\gets\beta_2\,Q_{jj}+(1-\beta_2)\,\|(G_A)_{:,j}\|_2^2$ \Comment{input side}
\State $P_{ii}\gets\beta_2\,P_{ii}+(1-\beta_2)\,\|(G_B)_{i,:}\|_2^2$ \Comment{output side}

\Phase{Whiten, polar, unwhiten}
\State $C_A\gets B^\top P B$, \qquad $C_B\gets A Q A^\top$ \Comment{small-side curvature, $r\times r$}
\State $W_A\gets C_A^{-1/2}\,\polar\!\big(C_A^{-1/2}\,\hat M_A\,Q^{-1/2}\big)\,Q^{-1/2}$
\State $W_B\gets P^{-1/2}\,\polar\!\big(P^{-1/2}\,\hat M_B\,C_B^{-1/2}\big)\,C_B^{-1/2}$

\Phase{Operator-norm step}
\State $\rho\gets\eta/(\smax(A)+\smax(B))$ \Comment{\Cref{eq:ours-prodcap}}
\State $\dot A\gets-\rho\,W_A/\smax(W_A)$, \qquad $A\gets A+\dot A$
\State $\dot B\gets-\rho\,W_B/\smax(W_B)$, \qquad $B\gets B+\dot B$
\end{algorithmic}
\end{algorithm}
\noindent Every $\smax(\cdot)$ is the guarded estimate of \Cref{alg:smax}, and every inverse-sqrt is relative-damped (\Cref{app:damping}).

\paragraph{Cost.} Beyond AdamW's two factor momenta:
\begin{itemize}[leftmargin=1.5em,itemsep=1pt,topsep=2pt]
  \item Extra memory $O(d_{\mathrm{in}}+d_{\mathrm{out}}+r^2)$: the two diagonals plus the two $r\times r$ eigenbases of $C_A,C_B$; no $O(d_{\mathrm{in}}d_{\mathrm{out}})$ dense state.
  \item Extra compute per step $O\!\big((d_{\mathrm{in}}+d_{\mathrm{out}})\,r^2\big)$ (one full polar per factor --- PolarExpress \cite{polarexpress} in Gram form \cite{gram_ns}), plus an eigenbasis refresh (QR with power-subspace iteration \cite{soap}) every $K{=}10$ steps, amortizing to $O(r^3/K)$.
  \item Measured wall overhead $1.01$--$1.06\times$ AdamW at production batch (\Cref{sec:exp-walltime}).
\end{itemize}

% ===========================================================================
\section{Experiments}
\label{sec:exp}

\subsection{Main result: speedup across 14 settings}
% SKELETON -- Figure 1 = AlgoPerf-style performance profile, \methodname vs AdamW
% vs the expensive spectral-polar members over all 14 workloads. The "we win" figure.
\begin{figure}[t]
  \centering
  % \includegraphics[width=\linewidth]{performance_profile.pdf}
  \fbox{\parbox{0.9\linewidth}{\centering\vspace{3em}
    \textbf{Figure 1 (placeholder).} Performance profile across 14 settings:
    fraction of settings where each method is within $\tau$ of the fastest.
    \methodname dominates. Generated after the coverage fill.\vspace{3em}}}
  \caption{Placeholder.}
  \label{fig:profile}
\end{figure}

\subsection{Ablations: what carries the speedup}
% SKELETON -- keystone figure. Knobs at 2-3 anchor cells:
%   (a) +polar vs -polar (curvature-only)        -> polar is essential
%   (b) dense vs diagonal curvature              -> flavor doesn't matter
%   (c) single-pass vs refined inner iteration   -> refinement doesn't matter
% NOTE: the -polar arm is nearly unrun today (1/14) -- MUST run.

\subsection{From step speedup to wall-clock}
\label{sec:exp-walltime}
% MAIN POINT (author's): a step-count speedup only matters if it survives in
% wall-clock. The optimizer.step is a fixed cost (batch-independent), so its share
% shrinks as batch grows -> the overhead amortizes. Numbers below are MEASURED
% (Blackwell, OLMo-2-1B, all-linear, seq=2048, full polar PE=8 vs AdamW) -- these
% are final, not placeholders.
\begin{table}[h]\centering\small
\begin{tabular}{lccc}
\toprule
per-step wall $/$ AdamW & global batch 4 & 16 (production) & 64 \\
\midrule
$r=64$  & $1.12\times$ & $1.03\times$ & $1.01\times$ \\
$r=256$ & $1.21\times$ & $1.06\times$ & $1.01\times$ \\
\bottomrule
\end{tabular}
\caption{Measured per-step wall overhead of \methodname{} (full polar) over AdamW;
8B/$r{=}256\approx 1.04\times$ at production batch. opt.step is fixed, so overhead
amortizes with batch.}
\end{table}
% TODO author: wall-clock speedup = (step speedup) / (per-step ratio); only the tiny
% global-batch-4 regime, which we do not run, costs >10%.

\subsection{Task dependence of the speedup}
\label{sec:exp-taskdep}
% MAIN POINT (author's): the speedup is task-dependent -- it grows with how OOD the
% task is and with rank; it is smallest on in-distribution code at the larger models.
% Numbers (speedup-vs-AdamW @ best lr; \PLACEHOLDER -- partial-polar data, refresh
% after the full-polar sweep):
%   - hard / in-distribution code ~1.2x: Llama-3-8B/opc r256 = 1.24x;
%     Qwen2.5-1.5B/opc r256 = 1.20x.
%   - OOD higher: Qwen2.5-1.5B/Bengali r256 = 1.57x (SAME model & rank as opc 1.20x
%     -- the clean controlled pair: only the dataset changes).
%   - rank: OLMo-2-1B/OpenMath r64 -> r256 = 1.50x -> 1.71x.
% TODO author: figure -- speedup vs (OOD proxy, rank).

% ===========================================================================
\section{Related work}
\label{sec:related}
% SKELETON -- organize by methodological axis, NOT paper-by-paper:
%   - spectral / steepest-descent updates: Muon and variants.
%   - Shampoo / SOAP curvature.
%   - Riemannian-preconditioned LoRA -- the product-norm precedent.
%   - LoRA optimizer/init work (GaLore, LoRA+, nonzero-init).
%   - rank-aware lr rules \cite{mua} -- cite as related work only; no prediction.

% ===========================================================================
\section{Limitations}
\label{sec:limitations}
% SKELETON -- TODO author: single base-model family scale ceiling (<=8B);
% eval-loss primary, downstream pass@1 limited/absent; single/low seed count
% -> state sigma coverage.

% ===========================================================================
\section{Conclusion}
\label{sec:conclusion}
% SKELETON -- TODO author.

% ===========================================================================
\appendix
\section{Proofs}
\label{app:proofs}

\begin{proof}[Proof of \Cref{lem:lmo}]
Let $M=U\Sigma V^\top$. For any $U'$ with $\|U'\|_2\le1$,
\[
  \langle M,U'\rangle=\sum_i\sigma_i\,u_i^\top U' v_i \;\ge\; -\sum_i\sigma_i,
\]
since $|u_i^\top U' v_i|\le\|U'\|_2\le1$. Equality holds at $U'=-\polar(M)$, uniquely when $M$ has full rank with distinct singular values; the zero and repeated singular directions otherwise leave $U'$ free.
\end{proof}

\begin{proof}[Proof of \Cref{prop:imuon}]
The objective separates over the two factors:
\[
  \langle G,\ \dot B A+B\dot A\rangle=\langle G_A,\dot A\rangle+\langle G_B,\dot B\rangle .
\]
Take the $A$-block under its budget $\|B\dot A\|_2\le\eta$. Write $B=U_B S_B^{1/2}$ with $U_B^\top U_B=I$, so that $\|B\dot A\|_2=\|S_B^{1/2}\dot A\|_2$. The substitution $Z=S_B^{1/2}\dot A$ turns the block into
\[
  \min_{\|Z\|_2\le\eta}\ \langle S_B^{-1/2}G_A,\,Z\rangle,
\]
which \Cref{lem:lmo} solves at $Z=-\eta\,\polar(S_B^{-1/2}G_A)$. Undoing the substitution gives
\[
  \dot A=S_B^{-1/2}Z=-\eta\,S_B^{-1/2}\polar\!\big(S_B^{-1/2}G_A\big),
\]
and the $B$-block is symmetric. This recovers iMuon \cite{imuon}, the per-factor split of compositional Muon \cite{tilde_csd}.
\end{proof}

\begin{proof}[Proof of \Cref{cor:imuon}]
With $U_B=BS_B^{-1/2}$ (orthonormal columns),
\[
  B\dot A=-\eta\,U_B\,\polar(S_B^{-1/2}G_A).
\]
Left-multiplication by $U_B$ preserves the spectral norm, and $\polar(\cdot)$ has unit nonzero singular values, so $\|B\dot A\|_2=\eta$.
\end{proof}

\begin{proof}[Proof of \Cref{prop:ours}]
Write $C_A=B^\top P B$ and substitute $Y_A=C_A^{1/2}\dot A\,Q^{1/2}$, so $\dot A=C_A^{-1/2}Y_A Q^{-1/2}$. Since $P^{1/2}B\,C_A^{-1/2}$ has orthonormal columns (its Gram is $I$), left-multiplication preserves singular values:
\[
  \|B\dot A\|_{P,Q}=\|P^{1/2}B\dot A\,Q^{1/2}\|_2=\|Y_A\|_2 .
\]
The objective is plain in $Y_A$ as well: with $g_A=B^\top G$ and trace cyclicity,
\[
  \langle B\dot A,-G\rangle=\langle\dot A,-g_A\rangle=\langle Y_A,-H_A\rangle,\qquad H_A=C_A^{-1/2}g_A Q^{-1/2}.
\]
So the problem reduces to $\max_{\|Y_A\|_2\le\tau}\langle Y_A,-H_A\rangle$, solved by $Y_A=-\tau\,\polar(H_A)$ (\Cref{lem:lmo}); unwhitening,
\[
  \dot A=-\tau\,C_A^{-1/2}\,\polar\!\big(C_A^{-1/2}g_A Q^{-1/2}\big)\,Q^{-1/2}.
\]
The Frobenius ball $\|Y_A\|_F\le\tau$ instead gives $\dot A\propto C_A^{-1}g_A Q^{-1}$, recovering AdaPreLoRA.
\end{proof}

\section{Spectral-norm estimation}
\label{app:smax}
\Cref{alg:smax} estimates every $\smax(\cdot)$ in \Cref{alg:ours}. The floor at the larger row/column $\ell_2$ norm (both lower-bound $\smax$) is load-bearing: a cold or stale power iteration can under-estimate $\smax$, and an under-estimate at the rescale $W/\smax(W)$ inflates the step past $\rho$ and can diverge.

\begin{algorithm}[h]
\caption{Guarded spectral-norm estimate $\smax(M)$}
\label{alg:smax}
\begin{algorithmic}[1]
\Require matrix $M$; warm-start vector $v$ (cached across steps); iterations $K$
\Statex

\State $\ell\gets\max\!\big(\max_i\|M_{i,:}\|_2,\ \max_j\|M_{:,j}\|_2\big)$
\Statex

\For{$k=1,\dots,K$} \Comment{power iteration on $M^\top M$}
  \State $w\gets Mv$
  \State $v\gets M^\top w \,/\, \|M^\top w\|_2$
\EndFor
\Statex

\State \Return $\max\!\big(\|Mv\|_2,\ \ell\big)$ \Comment{$\ell$ floor}
\end{algorithmic}
\end{algorithm}

\section{Relative damping}
\label{app:damping}
Each curvature inverse-sqrt in \Cref{alg:ours} is relative-damped: for a curvature factor $X$ with top eigenvalue $\lambda_{\max}$,
\begin{equation*}
  X^{-1/2}\ \longmapsto\ \big(X/\lambda_{\max}+\delta I\big)^{-1/2},\qquad \delta>0 .
\end{equation*}
The normalization and the floor each follow from the step's structure.
\begin{itemize}[leftmargin=1.5em,itemsep=2pt,topsep=2pt]
  \item \emph{Normalize by $\lambda_{\max}$.} The whitened polar depends only on the curvature's \emph{shape}, not its scale: $\polar\!\big(C_A^{-1/2}\hat M_A Q^{-1/2}\big)$ is invariant under $X\mapsto cX$, and the operator-norm rescale $W_A/\smax(W_A)$ reabsorbs any global factor. The scale-free representative is the normalized spectrum $X/\lambda_{\max}\in(0,1]$.
  \item \emph{Floor by $\delta$.} On that spectrum $\delta$ is one conditioning knob: it caps the metric's condition number at ${\sim}1/\delta$, so a near-null curvature direction cannot amplify the step without bound.
  \item \emph{Warmup.} A still-zero factor ($\lambda_{\max}=0$) maps to $I$, so steps before the EMAs warm up are plain momentum.
\end{itemize}

\printbibliography
\end{document}
